\section*{Capítulo \ref{ch:g8:winds-global-circulation} --- Los vientos y la circulación general}

\begin{solution}{exo:g8:winds-global-circulation:1}
(a) El alisio del nordeste, del nordeste. (b) Los \omterm{def:g8:winds-global-circulation:global-winds}{vientos del oeste}, del
noroeste. (c) Los \omterm{def:g8:winds-global-circulation:global-winds}{vientos polares del este}, del este (del nordeste).
(d) El alisio del sudeste, del sudeste.
\end{solution}

\begin{solution}{exo:g8:winds-global-circulation:2}
(a) El aire cálido sube y el aire frío baja. (b) Baja. (c) De las altas
presiones hacia las bajas presiones. (d) Hacia la derecha.
\end{solution}

\begin{solution}{exo:g8:winds-global-circulation:3}
(a) Falso: una \omterm{def:g8:winds-global-circulation:breezes}{brisa marina} sopla del \omterm{def:g2:water-on-earth:ocean}{mar} hacia la tierra (de día).
(b) Verdadero. (c) Falso: soplan del este (del nordeste o del sudeste).
(d) Verdadero.
\end{solution}

\begin{solution}{exo:g8:winds-global-circulation:4}
El dibujo muestra un circuito: a ras de suelo, una flecha del \omterm{def:g2:water-on-earth:ocean}{mar} hacia la
tierra; sobre la tierra, una flecha que sube; en altura, una flecha que
vuelve de la tierra sobre el \omterm{def:g2:water-on-earth:ocean}{mar}; sobre el \omterm{def:g2:water-on-earth:ocean}{mar}, una flecha que baja. La A
se marca en el suelo bajo el aire que baja (sobre el \omterm{def:g2:water-on-earth:ocean}{mar}), y la B bajo el
aire que sube (sobre la tierra).
\end{solution}

\begin{solution}{exo:g8:winds-global-circulation:5}
(a) 1.41 veces más suelo. (b) 50~\%. (c) 26~\%. (d) Cerca de los \omterm{def:g3:continents-and-oceans:poles}{polos}, la
luz del Sol llega inclinada y se reparte sobre una gran superficie, así que
cada metro cuadrado recibe poca energía.
\end{solution}

\begin{solution}{exo:g8:winds-global-circulation:6}
(a) \omterm{def:g8:winds-global-circulation:global-winds}{Vientos polares del este} (de \ang{70} a \ang{60} N), \omterm{def:g8:winds-global-circulation:global-winds}{vientos del oeste}
(de \ang{60} a \ang{30} N), \omterm{def:g8:winds-global-circulation:global-winds}{alisios} del nordeste (de \ang{30} N al
\omterm{def:g3:continents-and-oceans:equator}{ecuador}), \omterm{def:g8:winds-global-circulation:global-winds}{alisios} del sudeste (del \omterm{def:g3:continents-and-oceans:equator}{ecuador} a \ang{30} S), \omterm{def:g8:winds-global-circulation:global-winds}{vientos del oeste}
(de \ang{30} a \ang{60} S), \omterm{def:g8:winds-global-circulation:global-winds}{vientos polares del este} (de \ang{60} a
\ang{70} S). (b) Cerca del \omterm{def:g3:continents-and-oceans:equator}{ecuador}, en las calmas ecuatoriales. (c) Hacia
los \ang{30} N y los \ang{30} S, bajo el aire que baja.
\end{solution}

\begin{solution}{exo:g8:winds-global-circulation:7}
(a) La tierra ($31 > 22$): el viento sopla del \omterm{def:g2:water-on-earth:ocean}{mar} hacia la tierra, una
\omterm{def:g8:winds-global-circulation:breezes}{brisa marina}. (b) El \omterm{def:g2:water-on-earth:ocean}{mar} ($21 > 17$): el viento sopla de la tierra hacia
el \omterm{def:g2:water-on-earth:ocean}{mar}, una \omterm{def:g8:winds-global-circulation:breezes}{brisa terrestre}. (c) La tierra: $31 - 17 = 14$ grados; el \omterm{def:g2:water-on-earth:ocean}{mar}: $22 -
21 = 1$ grado. (d) Los rayos del Sol solo calientan una capa fina de suelo,
que se calienta y se enfría deprisa; el \omterm{def:g2:water-on-earth:ocean}{mar} reparte su calor en una capa
gruesa de agua, así que su temperatura apenas cambia entre el día y la
noche.
\end{solution}

\begin{solution}{exo:g8:winds-global-circulation:8}
(a) Del norte. (b) Hacia la derecha. (c) Del nordeste (es el alisio del
nordeste). (d) Sin rotación soplaría del sur; la rotación lo desvía hacia
la izquierda; sopla del sudeste.
\end{solution}

\begin{solution}{exo:g8:winds-global-circulation:9}
Cerca del \omterm{def:g3:continents-and-oceans:equator}{ecuador}, el suelo y el \omterm{def:g2:water-on-earth:ocean}{mar} están muy \emph{calentados} por el
Sol; el aire cálido y húmedo \emph{sube}; al subir \emph{se enfría}, y su
vapor de agua se convierte en \emph{\omterm{def:g4:clouds-rain-snow:cloud}{nubes}}, que dan chubascos y \omterm{def:g5:natural-hazards:storm}{tormentas}.
Hacia los \ang{25} N, el aire que subió en el \omterm{def:g3:continents-and-oceans:equator}{ecuador} \emph{baja}; debajo
hay \emph{altas presiones} en el suelo; el aire que baja se calienta, sus \omterm{def:g4:clouds-rain-snow:cloud}{nubes} se disipan y la región se mantiene
\emph{seca}: un \omterm{def:g5:weather-and-climate:desert}{desierto}.
\end{solution}

\begin{solution}{exo:g8:winds-global-circulation:10}
(a) Los \omterm{def:g8:winds-global-circulation:global-winds}{vientos del oeste}. No: soplan del oeste, en contra de un barco que
navega hacia el oeste. (b) Para llegar a la franja de los \omterm{def:g8:winds-global-circulation:global-winds}{alisios}, al sur
de unos \ang{30} N. (c) El alisio del nordeste, entre el \omterm{def:g3:continents-and-oceans:equator}{ecuador} y unos
\ang{30} N. (d) Los \omterm{def:g8:winds-global-circulation:global-winds}{vientos del oeste}, al norte de unos \ang{30} N (hacia
los \ang{40} N, la \omterm{def:g6:latitude-longitude:latitude}{latitud} del puerto de origen).
\end{solution}

\begin{solution}{exo:g8:winds-global-circulation:11}
(a) $1 \div 0.5 = 2$: la tabla da 2. (b) $1 \div 0.866 \approx
1.1547$, es decir, \textbf{1.15}. (c) A \ang{30}: $1000 \div 1.15 \approx
\qty{870}{W}$ por cada metro cuadrado (exactamente, $1000 \times 0.866 =
\qty{866}{W}$). A \ang{60}: $1000 \div 2 = \qty{500}{W}$.
\end{solution}

\begin{solution}{exo:g8:winds-global-circulation:12}
(a) Tres circuitos: el aire sube en el \omterm{def:g3:continents-and-oceans:equator}{ecuador} y baja hacia los \ang{30};
sube hacia los \ang{60} y baja hacia los \ang{30}; sube hacia los \ang{60}
y baja en el \omterm{def:g3:continents-and-oceans:poles}{polo}. (b) B en el \omterm{def:g3:continents-and-oceans:equator}{ecuador} y hacia los \ang{60}; A hacia los
\ang{30} y en el \omterm{def:g3:continents-and-oceans:poles}{polo}. (c) De los \ang{30} al \omterm{def:g3:continents-and-oceans:equator}{ecuador}, el aire de
superficie va hacia el sur y se desvía hacia la derecha: alisio del
nordeste. De los \ang{30} a los \ang{60}, va hacia el norte y se desvía
hacia la derecha: viene del sudoeste. Del \omterm{def:g3:continents-and-oceans:poles}{polo} a los \ang{60}, va hacia el
sur, desviado hacia la derecha: viento polar del este. (d) En la célula del
medio, el aire de superficie se mueve hacia el \omterm{def:g3:continents-and-oceans:poles}{polo}, es decir, hacia el
norte; desviado hacia la derecha, se dirige hacia el nordeste: son los
\omterm{def:g8:winds-global-circulation:global-winds}{vientos del oeste}.
\end{solution}

\begin{solution}{pb:g8:winds-global-circulation:1}
\textbf{1.} $59.6 - 23.5 = 36.1$ grados de \omterm{def:g6:latitude-longitude:longitude}{longitud}.

\textbf{2.} $36.1 \times 108 \approx 3899$: unos \qty{3900}{km}.

\textbf{3.} Sí: $3900 - 3896 = 4$~km, una diferencia minúscula.

\textbf{4.} Sí: los dos están cerca de los \ang{14} N, entre el \omterm{def:g3:continents-and-oceans:equator}{ecuador} y
unos \ang{30} N, la franja del alisio del nordeste.

\textbf{5.} Del nordeste (las flechas apuntan hacia el sudoeste).

\textbf{6.} Hacia delante, y un poco de lado: el viento viene de detrás
del barco, ligeramente por su derecha, y lo empuja hacia el oeste.

\textbf{7.} Forma parte de la \omterm{def:g8:winds-global-circulation:circulation}{circulación general}: el aire no deja de
bajar hacia los \ang{30} N ni de subir cerca del \omterm{def:g3:continents-and-oceans:equator}{ecuador}, así que el aire
de superficie no deja de moverse hacia el \omterm{def:g3:continents-and-oceans:equator}{ecuador}, y la rotación de la
Tierra siempre lo desvía hacia la derecha.

\textbf{8.} El alisio soplaría en contra del barco. El capitán debería
navegar hacia el norte, más allá de unos \ang{30} N, para encontrar los
\omterm{def:g8:winds-global-circulation:global-winds}{vientos del oeste}.

\textbf{9.} $7 \times 24 = \qty{168}{km}$ al día.

\textbf{10.} $3 \times 168 = \qty{504}{km}$.

\textbf{11.} $3896 \div 168 \approx 23.2$ días.

\textbf{12.} Unos \textbf{23 días}: menos de un mes.
\end{solution}
